% \iffalse meta-comment
% vim: textwidth=75
% Copyright (C) 2025 by Oliver Irwin <oliver@irwin.sh>
% -------------------------------------------------------------------------
% This work may be distributed and/or modified under the
% conditions of the LaTeX Project Public License (LPPL), either
% version 1.3c of this license or (at your option) any later
% version.  The latest version of this license is in the file:
%
% http://www.latex-project.org/lppl.txt
%
% This work is "maintained" (as per LPPL maintenance status) by
% Oliver Irwin.
%
% This work consists of the file liquidtex.dtx, a Makefile, and the derived
% filebase liquidtex.sty.
%
%<*internal>
\iffalse
%</internal>
%<*readme>
| liquidtex: | LaTeX with integrated Liquid templating                        |
|------------|----------------------------------------------------------------|
| Author:    | Oliver Irwin                                                   |
| E-mail:    | oliver@irwin.sh                                                |
| License:   | Released under the LaTeX Project Public License v1.3c or later |
| See:       | http://www.latex-project.org/lppl.txt                          |

Running `make` generates the derived files README, liquidtex.pdf and liquidtex.sty.
Running `make inst` installs the files in the user's TeX tree.
Running `make install` installs the files in the local TeX tree.

# liquidtex

`liquidtex` is a LaTeX package that integrates **Liquid templates**
directly into the LaTeX compilation process. It allows you to write LaTeX
documents that include dynamic placeholders and conditional logic, rendered
automatically from structured data sources (e.g. YAML, JSON) at compile
time.

## Features

- Seamless integration of **Liquid templating** within LaTeX documents.
- Support for **YAML** data sources.
- Verbose **logging** and **error reporting** for debugging template
  rendering.
- Configurable **options** for data source paths, compilation behaviour,
  and output verbosity.
- Compatible with `.sty` or `.dtx` distribution formats.

## Use Cases

- Auto-generated CVs or reports based on structured data.
- Dynamic lecture notes, syllabi, or publication lists.
- Research documents where repetitive elements are driven by metadata.
%</readme>
%<*ctanomat>
% This is a description file for ctan-o-mat.
% It manages uploads of a package to
% CTAN -- the Comprehensive TeX Archive Network.
%
% The syntax is roughly oriented towards (La)TeX.
% Two form of the macros are used. The simple macros take one argument
% in braces. Here the argument may not contain embedded macros.
%
% The second form uses an environment enclosed in \begin{}/\end{}.
% In the long text fields logo macros can be used.
%
% You should enter your values between the begin and the end of the
% named type.
% -------------------------------------------------------------------------
% This field contains the CTAN id of the package.
% The value is optional.
% The value is restricted to 32 characters.
\pkg{liquidtex}
% -------------------------------------------------------------------------
% This field contains the version of the package.
% The value is optional.
% The value is restricted to 32 characters.
\version{0.1.0 2025-11-13}
% -------------------------------------------------------------------------
% This field contains the name of the author(s).
% The value is optional.
% The value is restricted to 128 characters.
\author{Oliver Irwin}
% -------------------------------------------------------------------------
% This field contains the email address of the uploader.
% The value is an email address.
% The value is restricted to 255 characters.
\email{oliver@irwin.sh}
% -------------------------------------------------------------------------
% This field contains the name of the uploader.
% The value is restricted to 255 characters.
\uploader{Oliver Irwin}
% -------------------------------------------------------------------------
% This field contains the directory on CTAN.
% The value is optional.
% The value is restricted to 255 characters.
\ctanPath{macros/generic}
% -------------------------------------------------------------------------
% This field contains the license.
% The value is optional.
% The value is restricted to 2048 characters.
% Multiple values are allowed.
\license{lppl1.3c}
% -------------------------------------------------------------------------
% This field contains the URL of the home page.
% The value is optional.
% The value is a URL.
% The value is restricted to 255 characters.
% Multiple values are allowed.
\home{https://liquidtex.irwin.sh}
% -------------------------------------------------------------------------
% This field contains the URL of the bug tracker.
% The value is optional.
% The value is a URL.
% The value is restricted to 255 characters.
% Multiple values are allowed.
%\bugtracker{https://github.com/olirwin/liquidtex/issues/}
% -------------------------------------------------------------------------
% This field contains the support channel.
% The value is optional.
% The value is a URL.
% The value is restricted to 255 characters.
% Multiple values are allowed.
\support{}
% -------------------------------------------------------------------------
% This field contains the version management.
% The value is optional.
% The value is a URL.
% The value is restricted to 255 characters.
% Multiple values are allowed.
%\repository{https://github.com/olirwin/liquidtex}
% -------------------------------------------------------------------------
% This field contains the developer's channel.
% The value is optional.
% The value is a URL.
% The value is restricted to 255 characters.
% Multiple values are allowed.
\development{}
% -------------------------------------------------------------------------
% This field contains the update indicator; true for update.
% The value is optional.
% The value is restricted to 8 characters.
\update{false}
% -------------------------------------------------------------------------
% This field contains the topics id.
% The value is optional.
% The value is restricted to 1024 characters.
% Multiple values are allowed.
\topic{}
% -------------------------------------------------------------------------
% This field contains the text for the mail announcement.
% The value is optional.
% The value is restricted to 8192 characters.
\begin{announcement}
\end{announcement}
% -------------------------------------------------------------------------
% This field contains the one-liner for the package.
% The value is optional.
% The value is restricted to 128 characters.
\summary{A LaTeX package to use integrated Liquid templates in documents}
% -------------------------------------------------------------------------
% This field contains the descriptive abstract for the package.
% The value is optional.
% The value is restricted to 4096 characters.
\begin{description}
# liquidtex

`liquidtex` is a LaTeX package that integrates **Liquid templates**
directly into the LaTeX compilation process. It allows you to write LaTeX
documents that include dynamic placeholders and conditional logic, rendered
automatically from structured data sources (e.g. YAML, JSON) at compile
time.

## Features

- Seamless integration of **Liquid templating** within LaTeX documents.
- Support for **YAML** data sources.
- Verbose **logging** and **error reporting** for debugging template
  rendering.
- Configurable **options** for data source paths, compilation behaviour,
  and output verbosity.
- Compatible with `.sty` or `.dtx` distribution formats.

## Use Cases

- Auto-generated CVs or reports based on structured data.
- Dynamic lecture notes, syllabi, or publication lists.
- Research documents where repetitive elements are driven by metadata.
\end{description}
% -------------------------------------------------------------------------
% This field contains the note to the CTAN upload managers.
% The value is optional.
% The value is restricted to 2048 characters.
\begin{note}
\end{note}
% -------------------------------------------------------------------------
% This field contains the archive file.
% The value is the file name of the archive to be uploaded.
% It may have a relative or absolute directory.
\file{liquidtex.zip}
%</ctanomat>
%<*changelog>
%</changelog>
%
%<*internal>
\fi
\def\nameofplainTeX{plain}
\ifx\fmtname\nameofplainTeX\else
  \expandafter\begingroup
\fi
%</internal>
%<*package|install>
\def\liquidtexversion{0.1.0}
\def\liquidtexdate{2025-11-12}
%</package|install>
%<*install>
\input l3docstrip.tex
\Msg{********************************************************}
\Msg{* Installation}
\Msg{* Package: liquidtex \liquidtexdate\space v\liquidtexversion}
\Msg{* that allows for integrated Liquid templating in documents}
\Msg{********************************************************}
\keepsilent
\askforoverwritefalse
\preamble
---------:| ---------------------------------------------------------------
liquidtex:| LaTeX with integrated Liquid templating
   Author:| Oliver Irwin
   E-mail:| oliver@irwin.sh
  License:| Released under the LaTeX Project Public License v1.3c or later
      See:| http://www.latex-project.org/lppl.txt

\endpreamble
\declarepreamble\utilpreamble
\endpreamble
\postamble

Copyright (C) 2025 by Oliver Irwin <oliver@irwin.sh>

This work may be distributed and/or modified under the
conditions of the LaTeX Project Public License (LPPL), either
version 1.3c of this license or (at your option) any later
version.  The latest version of this license is in the file:

http://www.latex-project.org/lppl.txt

This work is "maintained" (as per LPPL maintenance status) by
Oliver Irwin.

This work consists of the file `liquidtex.dtx`, the .sty files located in the
`utils` folder and a Makefile.
Running "make" generates the derived files README, liquidtex.pdf and liquidtex.sty.
Running "make inst" installs the files in the user's TeX tree.
Running "make install" installs the files in the local TeX tree.

\endpostamble
\declarepostamble\utilpostamble
\endpostamble
%\usedir{tex/latex/liquidtex}
\def\LQT@pkgname{\jobname}
\def\LQT@packagename{\LQT@pkgname doc}
%</install>
%<*install>
%\usedir{source/latex/liquidtex}
\usedir{.}
\generate{
  \nopreamble\nopostamble
  \file{\LQT@pkgname.pkg}{\from{\jobname.dtx}{ctanomat}}
}%
\usedir{doc/latex/\LQT@pkgname}
\generate{
  \nopreamble\nopostamble
  \file{README.md}{\from{\jobname.dtx}{readme}}
}%
\generate{
  \nopreamble\nopostamble
  \file{CHANGELOG.md}{\from{\jobname.dtx}{changelog}}
}%
\usedir{tex/latex/\LQT@pkgname}
\generate{
  \file{\jobname.sty}{%
    \from{\jobname.dtx}{package}%
  }
  \file{lqt-messages.sty}{%
    \usepreamble\utilpreamble\usepostamble\utilpostamble
    \from{utils/lqt-messages.dtx}{preamble}%
    \from{utils/lqt-messages.dtx}{package}%
    \from{utils/lqt-messages.dtx}{postamble}%
  }
  \file{lqt-files.sty}{%
    \usepreamble\utilpreamble\usepostamble\utilpostamble
    \from{utils/lqt-files.dtx}{preamble}%
    \from{utils/lqt-files.dtx}{package}%
    \from{utils/lqt-files.dtx}{postamble}%
  }
  \file{lqt-engine-external.sty}{%
    \usepreamble\utilpreamble\usepostamble\utilpostamble
    \from{utils/lqt-engine-external.dtx}{preamble}%
    \from{utils/lqt-engine-external.dtx}{package}%
    \from{utils/lqt-engine-external.dtx}{postamble}%
  }
}
%</install>
%<install>\endbatchfile
%<*internal>
\usedir{source/latex/\LQT@pkgname}
\generate{
  \file{\LQT@pkgname.ins}{\from{\jobname.dtx}{install}}
  \file{\LQT@pkgname.drv}{\from{\jobname.dtx}{driver}}%
}%
\ifx\fmtname\nameofplainTeX
  \expandafter\endbatchfile
\else
  \expandafter\endgroup
\fi
%</internal>
% \fi
%
% \iffalse
%<*driver>
\ProvidesFile{liquidtex.dtx}
%</driver>
%<*driver>
\documentclass[full,british]{l3doc}
\usepackage[a4paper,margin=25mm,left=50mm,nohead]{geometry}
\usepackage[yad]{denisbdoc}
\OnlyDescription
\DisableImplementation
\EnableCrossrefs
% \CodelineIndex
\RecordChanges
\begin{document}
  \DocInput{\jobname.dtx}
\end{document}
%</driver>
% \fi
% \CheckSum{4290}
%
% \CharacterTable
%  {Upper-case    \A\B\C\D\E\F\G\H\I\J\K\L\M\N\O\P\Q\R\S\T\U\V\W\X\Y\Z
%   Lower-case    \a\b\c\d\e\f\g\h\i\j\k\l\m\n\o\p\q\r\s\t\u\v\w\x\y\z
%   Digits        \0\1\2\3\4\5\6\7\8\9
%   Exclamation   \!     Double quote  \"     Hash (number) \#
%   Dollar        \$     Percent       \%     Ampersand     \&
%   Acute accent  \'     Left paren    \(     Right paren   \)
%   Asterisk      \*     Plus          \+     Comma         \,
%   Minus         \-     Point         \.     Solidus       \/
%   Colon         \:     Semicolon     \;     Less than     \<
%   Equals        \=     Greater than  \>     Question mark \?
%   Commercial at \@     Left bracket  \[     Backslash     \\
%   Right bracket \]     Circumflex    \^     Underscore    \_
%   Grave accent  \`     Left brace    \{     Vertical bar  \|
%   Right brace   \}     Tilde         \~}
%
%
%   \GetFileInfo{\jobname.dtx}
%   \DoNotIndex{%
%   \newcommand,\newenvironment,\begin,\end,\acrpluralsuffix,\acrshort,\active,\addlinespace,\advance,\AtEndDocument,\begingroup,\caption,\catcode,\ClassInfo,\color,\colorlet,\cr,\DeclareGraphicsExtensions,\declaretheorem,\declaretheoremstyle,\def,\definecolor,\do,\edef,\else,\emph,\endgroup,\endVerbatimOut,\expandafter,\fi,\floatsetup,\footnotesize,\gls,\glsentrylong,\glspluralsuffix,\glsshorttok,\graphicspath,\hide,\ifluatex,\includegraphics,\InputIfFileExists,\intertitle,\killienc,\label,\LaTeX,\linewidth,\LoadClass,\loadglsentries,\lstinputlisting,\lstMakeShortInline,\lstset,\LTXtable,\makeglossaries,\mdfdefinestyle,\mdfsetup,\medskipamount,\meta,\multicolumn,\noexpand,\normalfont,\nothing,\PassOptionsToClass,\PassOptionsToPackage,\ProcessOptions,\protect,\raisebox,\relax,\renewcommand,\RequirePackage,\temp,\TeX,\textbf,\textlangle,\textrangle,\textsc,\textsf,\texttt,\textup,\the,\tmp,\to,\todoformat,\todos,\ttdefault,\ttfamily,\typeout,\uccode,\uppercase,\VerbatimOut,\WithSuffix,\LoadClassWithOptions,\DeclareOption,\xspace,\undefined,\newif,\nothtfalse,\nothttrue,\HCode,\ifnotht,\ifx,\@bsphack}
%   \DoNotIndex{%
%   \@documentclasshook,\@empty,\@esphack,\@tempa,\\,\actualchar,\baselineskip,\boolean,\c@CodelineNo,\c@HD@hypercount,\count@,\Describe@Option,\Describe@ShortCut,\encapchar,\endoption,\endshortcut,\endtrivlist,\equal,\fvset,\g@addto@macro,\global,\hbox,\HD@target,\HDorg@encapchar,\hologo,\if@inlabel,\ignorespaces,\index,\input,\item,\leavevmode,\let,\levelchar,\llap,\long,\loop,\Lsmcp,\lstdefinestyle,\lstnewenvironment,\m@cro@YT,\m@ne,\macro@cnt,\MacroTopsep,\makelabel,\MakePrivateLetters,\marginnote,\marginpar,\newboolean,\nobreak,\or,\PackageWarning,\PrintDescribeEnv,\PrintDescribeOption,\PrintDescribeShortCut,\PrintEnvName,\PrintOptionName,\PrintShortCutName,\ProvidesFile,\raggedleft,\repeat,\reversedvideodbend,\saved@macroname,\setboolean,\sffamily,\shortcut,\space,\special@index,\SpecialMainOptionIndex,\SpecialMainShortCutIndex,\SpecialOptionIndex,\SpecialShortCutIndex,\string,\strut,\textbelow,\textsl,\topsep,\trivlist,\VerbatimEnvironment,\vss,\vtop,\xspaceaddexceptions,\par}
%
% \DoNotIndex{\@cclvi,\@ne,\@tempcnta,\@whilenum,\toks@,\z@}
%
% \DoNotIndex{\addbibresource,\biolinumKeyGlyph,\DeclareRobustCommand,\ClassWarning,\DescribeOption,\DescribeShortCut,\ifcase,\ifnum,\ifthenelse,\option,\pagestyle}
% \DoNotIndex{\newcommand,\newenvironment}
%
% \def\liquidtexversion{0.1.0}
% \def\liquidtexdate{2025-11-13}
%
%\title{\textsf{liquidtex} --- \LaTeX{} with integrated Liquid templating\thanks{This file
%   describes version \liquidtexversion, last revised \liquidtexdate.}
%}
%\author{Oliver Irwin\thanks{E-mail: oliver@irwin.sh}}
%\date{\liquidtexdate}
%
%\maketitle
%
% \begin{abstract}
%   The \pkg{liquidtex} package integrates a Liquid-style templating layer
%   into the \LaTeX{} compilation process.  It allows authors to incorporate
%   templated, data-driven content into their documents by automatically
%   rendering Liquid templates at compile time.
% \end{abstract}
%
% \begin{documentation}
%
% \section{Package options}
%
% The package uses a key--value interface based on \pkg{l3keys2e}.  Options
% may be passed at package load time or via \cs{liquidsetup} in the
% preamble.  Some keys are restricted to preamble usage only.
%
% \subsection{General options}
%
% \begin{description}
%
%   \item[\texttt{verbose} (bool)]
%     Enables verbose diagnostic messages.\newline
%     Default: \texttt{false}.
%
%   \item[\texttt{no-cache} (bool)]
%     Disables the use of cached rendered output.\newline
%     Default: \texttt{false}.
%
%   \item[\texttt{outdir = \meta{path}}]
%     Directory used to store rendered templates and cache files.
%     Default: \texttt{.liquid-cache}.\newline
%     This key may only be used in the preamble.
%
%   \item[\texttt{use-auxdir} (bool)]
%     Indicates whether \texttt{outdir} should be interpreted relative to
%     the engine’s auxiliary directory (e.g.\ under \texttt{--auxdir} or
%     \texttt{--output-directory}).  This integration is not yet
%     implemented; the option is reserved for future versions.\newline
%     Default: \texttt{true}.  Preamble-only.
%
% \end{description}
%
% \subsection{Template data}
%
% \begin{description}
%
%   \item[\texttt{data = \meta{file}}]
%     Sets the default data source used when rendering templates.
%     The value may be overridden for individual calls of
%     \cs{inputliquid}.  No default is imposed.
%
% \end{description}
%
% \subsection{Engine selection}
%
% Liquid templates can be rendered by different backends.  The engine is
% selected with:
%
% \begin{description}
%
%   \item[\texttt{engine = lua \textbar{} external}]
%     Chooses the rendering backend.  This option is restricted to the
%     preamble.
%
%     \begin{description}
%
%       \item[lua]
%         Uses a Lua-based renderer.  Requires compilation with Lua\LaTeX.
%
%       \item[external]
%         Uses an external program to render the template.  This backend
%         requires additional configuration through \texttt{engine-conf}.
%         Shell escape is required to run the external renderer.
%
%     \end{description}
%
% \end{description}
%
% \subsection{External engine configuration}
%
% When \texttt{engine = external}, the external renderer must be configured
% via:
%
% \begin{description}
%
%   \item[\texttt{engine-conf = \meta{key--value list}}]
%     A subkey list describing the external renderer, for example:
%
% \begin{verbatim}
%   engine-conf = {
%     path           = /usr/bin/liquidtex-cli,
%     arg-style      = positional,
%     template-index = 1,
%     data-index     = 2,
%     output-index   = 3,
%   }
% \end{verbatim}
%
% Supported sub-keys:
%
%     \begin{description}
%
%       \item[\texttt{path = \meta{string}}]
%         Absolute or relative path to the external executable.
%         This is mandatory when using the external engine.
%
%       \item[\texttt{arg-style = positional \textbar{} named}]
%         Determines how command-line arguments are passed to the external
%         renderer.
%
%       \item[\texttt{positional \textbar{} named}]
%         As a (maybe temporary) shortcut for the \texttt{arg-style} key options.
%
%       \item[\texttt{template-index}, \texttt{data-index}, \texttt{output-index}]
%         When using positional argument style, these keys specify the
%         1-based positions of the template, data, and output file arguments.
%         Defaults: \texttt{1}, \texttt{2}, \texttt{3}.
%
%     \end{description}
%
% This option is restricted to the preamble.  If \texttt{engine = external}
% and no \texttt{path} is provided, the package raises an error at
% \cs{begin}\env{\{document\}}.
%
% \end{description}
%
%
% \section{User interface}
%
% \subsection{Runtime configuration}
%
% \begin{function}{\liquidsetup}
% \begin{syntax}
%   \cs{liquidsetup}\marg{options}
% \end{syntax}
%
% Applies configuration keys at runtime.  The effect is identical to using
% package options, except that keys marked as “preamble-only” may only be
% used before \cs{begin}\env{\{document\}}.  Example:
%
% \begin{verbatim}
%   \liquidsetup{
%     outdir     = .liquid-cache,
%     use-auxdir = false,
%     verbose    = true,
%     data       = info.yml
%   }
% \end{verbatim}
%
% \end{function}
%
% \subsection{Cache directory management}
%
% The cache directory stores all rendered template output files.
%
% \begin{itemize}
%   \item The directory is determined by \texttt{outdir}.
%   \item It is created automatically at \cs{begin}\env{\{document\}}.
%   \item This requires shell escape, as directory creation uses
%         \cs{sys_shell_now:n}.  If shell escape is disabled, the package
%         issues an error asking the user to create the directory manually.
%   \item When \texttt{no-cache} is set, cache lookup is disabled and
%         templates are always re-rendered.
% \end{itemize}
%
% Throughout the compilation, the package may emit diagnostic messages
% about cache usage when \texttt{verbose} is enabled.
%
% \subsection{Liquid template rendering}
%
% \begin{function}{\inputliquid}
% \begin{syntax}
%   \cs{inputliquid}\oarg{data file}\marg{template file}
% \end{syntax}
%
% Renders a Liquid template and inputs the resulting \LaTeX{} code.
%
% \paragraph{Arguments}
% \begin{itemize}
%   \item \meta{template file}: mandatory path to a Liquid template.
%   \item \meta{data file}: optional override of the global \texttt{data}
%         setting.
% \end{itemize}
%
% \paragraph{Caching}
%
% For each pair \marg{template}\marg{data}, the package generates a cache
% file whose name is derived from the base names of the sources:
%
% \begin{quote}
%   \verb|<outdir>/<template>+<data>.tex|
% \end{quote}
%
% If no data file is used, the name simplifies to:
%
% \begin{quote}
%   \verb|<outdir>/<template>.tex|
% \end{quote}
%
% A cache entry is considered \emph{fresh} when:
%
% \begin{itemize}
%   \item it exists, and
%   \item it is not older than the template file, and
%   \item if a data file is present, it is not older than that data file.
% \end{itemize}
%
% Freshness is checked using the \pkg{filemod} package.
% If the cache is fresh, the cached file is input directly; otherwise, the
% rendering engine is invoked and the result stored.
%
% \paragraph{Rendering}
%
% Currently, \cs{inputliquid} is implemented only for the
% \texttt{external} engine.  Rendering proceeds as follows:
%
% \begin{itemize}
%   \item The external executable is invoked through the system shell.
%   \item Shell escape must therefore be enabled.
%   \item The template, data file, and output path are passed according to
%         the configuration set by \texttt{engine-conf}.
%   \item Upon success, the generated file is stored in the cache
%         directory and immediately input.
% \end{itemize}
%
% \paragraph{Example}
%
% \begin{verbatim}
%   \usepackage[
%     engine      = external,
%     engine-conf = { path = liquidtex-cli },
%     data        = cv.yml,
%     outdir      = .liquid-cache
%   ]{liquidtex}
%
%   % Uses cv.yml
%   \inputliquid{templates/header.liq}
%
%   % Override data for this template
%   \inputliquid[teaching.yml]{templates/teaching-entry.liq}
% \end{verbatim}
%
% \end{function}
%
% \end{documentation}
% \MaybeStop{
% }
% \begin{implementation}
%
% \subsection{Package declaration and internal variables}
%
%    \begin{macrocode}
%<*package>
\RequirePackage{expl3,xparse}
\ProvidesExplPackage
  {liquidtex}
  {\liquidtexdate}
  {\liquidtexversion}
  {LaTeX with integrated Liquid templating}

\ExplSyntaxOn

% internal configuration variables
\bool_new:N \__lqt_verbose_bool
\bool_new:N \__lqt_use_auxdir_bool
\bool_new:N \__lqt_no_cache_bool
\tl_new:N   \__lqt_engine_tl
\tl_new:N   \__lqt_outdir_tl
\tl_new:N   \__lqt_data_path_tl

\tl_new:N   \__lqt_engine_path_tl
\tl_new:N   \__lqt_engine_arg_style_tl
\int_new:N  \__lqt_engine_template_index_int
\int_new:N  \__lqt_engine_data_index_int
\int_new:N  \__lqt_engine_output_index_int

% internal derived value: effective cache directory
\tl_new:N   \__lqt_cache_dir_tl

% temp variables
\tl_new:N   \l__lqt_output_file_tl
\tl_new:N   \l__lqt_cmd_line_tl
\tl_new:N   \l__lqt_tmpl_base_tl
\tl_new:N   \l__lqt_data_base_tl
\tl_new:N   \l__lqt_tmp_dir_tl
\tl_new:N   \l__lqt_tmp_ext_tl
\tl_new:N   \l__lqt_tmp_template_tl
\tl_new:N   \l__lqt_tmp_data_tl
\tl_new:N   \l__lqt_engine_log_file_tl

%    \end{macrocode}
%
% \subsection{Key definitions}
%
% We define a key family \texttt{liquidtex} with one boolean key.
%
%    \begin{macrocode}
\keys_define:nn { liquidtex }
  {
    verbose .bool_set:N = \__lqt_verbose_bool,
    verbose .initial:n  = false,
    verbose .default:n  = true,

    outdir  .tl_set:N   = \__lqt_outdir_tl,
    outdir  .initial:n  = {.liquid-cache},
    outdir  .usage:n    = preamble,

    use-auxdir .bool_set:N = \__lqt_use_auxdir_bool,
    use-auxdir .initial:n  = true,
    use-auxdir .default:n  = true,
    use-auxdir .usage:n    = preamble,

    no-cache .bool_set:N = \__lqt_no_cache_bool,
    no-cache .initial:n  = false,
    no-cache .default:n  = true,
    no-cache .usage:n    = preamble,

    engine .choice:,
    engine / lua .code:n      =
    { \tl_set:Nn \__lqt_engine_tl { lua } },
    engine / external .code:n =
    { \tl_set:Nn \__lqt_engine_tl { external } },
    engine / unknown .code:n =
    {
      \msg_error:nnx { liquidtex } { unknown-engine } {#1}
    },

    engine .initial:n = external,
    engine .usage:n   = preamble,

    engine-conf .code:n =
    { \keys_set:nn { liquidtex / engine-conf } { #1 } },
    engine-conf .usage:n = preamble,

    data .tl_set:N   = \__lqt_data_path_tl,
    data .initial:n  = {},
  }

\keys_define:nn { liquidtex / engine-conf }
{
  path .tl_set:N   = \__lqt_engine_path_tl,
  path .initial:n  = {},

  positional .code:n = { \tl_set:Nn \__lqt_engine_arg_style_tl { positional } },
  named      .code:n = { \tl_set:Nn \__lqt_engine_arg_style_tl { named } },

  arg-style .choice:,
  arg-style / positional .code:n =
  { \tl_set:Nn \__lqt_engine_arg_style_tl { positional } },
  arg-style / named .code:n =
  { \tl_set:Nn \__lqt_engine_arg_style_tl { named } },
  arg-style .initial:n = positional,

  template-index .int_set:N = \__lqt_engine_template_index_int,
  template-index .initial:n = 1,

  data-index .int_set:N = \__lqt_engine_data_index_int,
  data-index .initial:n = 2,

  output-index .int_set:N = \__lqt_engine_output_index_int,
  output-index .initial:n = 3,
}
%    \end{macrocode}
% Process package options
%
%    \begin{macrocode}
\ProcessKeyOptions[liquidtex]
%    \end{macrocode}
%
% \subsection{Package requirements}
%    \begin{macrocode}
%% useful for file modification information - notably cache here
\RequirePackage{filemod}
%% internal message module
\RequirePackage{lqt-messages}
\RequirePackage{lqt-files}
\RequirePackage{lqt-engine-external}
%    \end{macrocode}
% \subsection{Public configuration command}
%
%    \begin{macrocode}
\NewDocumentCommand{\liquidsetup}{ m }{ \keys_set:nn { liquidtex } { #1 } }
%    \end{macrocode}
%
% \subsection{Liquid rendering}
%    \begin{macrocode}


\cs_new_protected:Npn \__lqt_input_liquid:nn #1#2
  {
    % Check that the template file exists
    \IfFileExists{ #1 }
      {}
      {
        \msg_error:nnx
          { liquidtex } { missing-template }
          { #1 }
      }

    \tl_if_eq:VnTF \__lqt_engine_tl { external }
      { \__lqt_render:nn { #1 }{ #2 } }
      {
        \msg_error:nnx
          { liquidtex } { engine-not-external }
          { \tl_use:N \__lqt_engine_tl }
      }
  }

\NewDocumentCommand \inputliquid { O{} m }
  {
    % #1 = optional data override, #2 = template
    \tl_set:Nn \l__lqt_tmp_template_tl { #2 }
    \tl_set:Nx \l__lqt_tmp_data_tl
      {
        \tl_if_blank:nTF { #1 }
          { \__lqt_data_path_tl }
          { #1 }
      }
    \__lqt_input_liquid:nn
      { \tl_use:N \l__lqt_tmp_template_tl }
      { \tl_use:N \l__lqt_tmp_data_tl }
  }
  
%    \end{macrocode}
%    \begin{macrocode}
\AtBeginDocument{ %
  \__lqt_init_cache_dir:
}
\ExplSyntaxOff
%</package>
%    \end{macrocode}
%
% \end{implementation}
%\Finale
\endinput
% Local Variables:
% mode: docTeX
% TeX-master: t
% End:
